% ============================================================
\chapter*{Chapitre 0 --- Installation et Configuration}
\addcontentsline{toc}{chapter}{Chapitre 0 --- Installation et Configuration}
\label{ch:installation}
% ============================================================

Avant de commencer \`a programmer, il faut pr\'eparer votre environnement de travail. Ce chapitre vous guide pas \`a pas.

\section*{Installer Python}
\addcontentsline{toc}{section}{Installer Python}

\begin{intuition}
Python est un \textbf{langage de programmation} gratuit, utilis\'e par des millions de scientifiques dans le monde. C'est comme une calculatrice surpuissante qui comprend vos instructions \'ecrites en anglais.
\end{intuition}

\subsection*{M\'ethode recommand\'ee : Anaconda}

\textbf{Anaconda}\index{Anaconda} est une distribution Python qui inclut toutes les biblioth\`eques scientifiques dont nous aurons besoin.

\begin{enumerate}
\item Rendez-vous sur \url{https://www.anaconda.com/download}
\item T\'el\'echargez la version pour votre syst\`eme (Windows, macOS ou Linux).
\item Lancez l'installateur et suivez les instructions par d\'efaut.
\item V\'erifiez l'installation en ouvrant un terminal et en tapant :
\end{enumerate}

\begin{codeblock}[title=\textbf{Terminal}]
\begin{minted}{bash}
python --version
\end{minted}
\end{codeblock}

\begin{codeoutput}
\begin{verbatim}
Python 3.12.4
\end{verbatim}
\end{codeoutput}

\subsection*{Alternative l\'eg\`ere : Python seul + pip}

Si vous pr\'ef\'erez une installation minimale :
\begin{enumerate}
\item T\'el\'echargez Python depuis \url{https://www.python.org/downloads/}
\item Cochez <<~Add Python to PATH~>> lors de l'installation.
\item Installez les biblioth\`eques n\'ecessaires :
\end{enumerate}

\begin{codeblock}[title=\textbf{Terminal}]
\begin{minted}{bash}
pip install numpy matplotlib pandas scipy seaborn scikit-learn jupyter
\end{minted}
\end{codeblock}

\section*{Jupyter Notebook}
\addcontentsline{toc}{section}{Jupyter Notebook}

\textbf{Jupyter Notebook}\index{Jupyter} est un environnement interactif o\`u vous \'ecrivez du code, voyez les r\'esultats et ajoutez des notes --- comme un cahier de laboratoire num\'erique.

\begin{codeblock}[title=\textbf{Terminal --- Lancer Jupyter}]
\begin{minted}{bash}
jupyter notebook
\end{minted}
\end{codeblock}

Votre navigateur s'ouvre automatiquement. Cliquez sur \textbf{New $\to$ Python 3} pour cr\'eer un nouveau notebook.

\begin{center}
\begin{tikzpicture}[
  box/.style={rectangle, draw, rounded corners, fill=#1, minimum width=4cm, minimum height=1cm, align=center, font=\small}
]
\node[box=blue!10] (cell1) at (0,2) {\texttt{In [1]: 2 + 3}};
\node[box=green!10] (out1) at (0,0.8) {\texttt{Out[1]: 5}};
\node[box=blue!10] (cell2) at (0,-0.5) {\texttt{In [2]: print("Bonjour")}};
\node[box=green!10] (out2) at (0,-1.7) {\texttt{Bonjour}};
\draw[-{Stealth},thick] (cell1.south) -- (out1.north);
\draw[-{Stealth},thick] (cell2.south) -- (out2.north);
\node[font=\footnotesize\itshape, right] at (3,2) {Cellule de code};
\node[font=\footnotesize\itshape, right] at (3,0.8) {R\'esultat};
\end{tikzpicture}
\end{center}

\begin{bonnepratique}
Utilisez \textbf{Shift + Entr\'ee} pour ex\'ecuter une cellule et passer \`a la suivante. Utilisez \textbf{Ctrl + Entr\'ee} pour ex\'ecuter sans changer de cellule.
\end{bonnepratique}

\section*{\'Editeurs de texte}
\addcontentsline{toc}{section}{\'Editeurs de texte}

Pour \'ecrire des scripts Python (fichiers \texttt{.py}), vous pouvez utiliser :
\begin{itemize}
\item \textbf{VS Code}\index{VS Code} (recommand\'e) : gratuit, avec extension Python.
\item \textbf{Spyder} : inclus dans Anaconda, con\c{c}u pour les scientifiques.
\item \textbf{PyCharm} : version Community gratuite.
\end{itemize}

\section*{V\'erification de l'installation}
\addcontentsline{toc}{section}{V\'erification de l'installation}

\begin{codeblock}[title=\textbf{Python --- Premier test}]
\begin{minted}{python}
import numpy as np
import matplotlib.pyplot as plt

print("Installation réussie !")
print(f"NumPy version : {np.__version__}")

# Petit test : tracer sin(x)
x = np.linspace(0, 2 * np.pi, 100)
plt.plot(x, np.sin(x))
plt.title("sin(x) — Votre premier graphique !")
plt.xlabel("x")
plt.ylabel("sin(x)")
plt.grid(True)
plt.savefig('test_installation.pdf')
plt.show()
print("Graphique sauvegardé !")
\end{minted}
\end{codeblock}

\begin{codeoutput}
\begin{verbatim}
Installation réussie !
NumPy version : 1.26.4
Graphique sauvegardé !
\end{verbatim}
\end{codeoutput}

Si tout fonctionne, vous \^etes pr\^et(e) \`a commencer le chapitre~1 !

\section*{Installer R (pour le chapitre 9)}
\addcontentsline{toc}{section}{Installer R}

\begin{enumerate}
\item T\'el\'echargez R depuis \url{https://cran.r-project.org/}
\item Installez \textbf{RStudio Desktop} depuis \url{https://posit.co/download/rstudio-desktop/}
\item V\'erifiez en ouvrant RStudio et en tapant :
\end{enumerate}

\begin{codeblock}[title=\textbf{Console R}]
\begin{minted}{r}
R.version.string
\end{minted}
\end{codeblock}

\begin{codeoutput}
\begin{verbatim}
[1] "R version 4.4.1 (2024-06-14)"
\end{verbatim}
\end{codeoutput}

\begin{errmark}
\begin{itemize}
\item \textbf{<<~Python n'est pas reconnu~>>} : vous avez oubli\'e d'ajouter Python au PATH. R\'einstallez en cochant l'option.
\item \textbf{<<~ModuleNotFoundError~>>} : la biblioth\`eque n'est pas install\'ee. Tapez \texttt{pip install nom\_du\_module}.
\item \textbf{Permission refus\'ee} : sur macOS/Linux, essayez \texttt{pip install --user nom\_du\_module}.
\end{itemize}
\end{errmark}
